%%%%%%%%%%%%%%%%%%%%%%%%%%%%%%%%%%%%%%%%%%%%%%%%%%%%%%%%%%
%%
%% MARK: Title
%%
%%%%%%%%%%%%%%%%%%%%%%%%%%%%%%%%%%%%%%%%%%%%%%%%%%%%%%%%%%

\chapter*{\S B17. Differential Forms on the Cross}
\setcounter{footnote}{0}

%%%%%%%%%%%%%%%%%%%%%%%%%%%%%%%%%%%%%%%%%%%%%%%%%%%%%%%%%%
%%
%% MARK: Abstract
%%
%%%%%%%%%%%%%%%%%%%%%%%%%%%%%%%%%%%%%%%%%%%%%%%%%%%%%%%%%%

\begin{abstract}
  We describe the differential forms on the diffeological space consisting of the union of the two axes $ox \cup oy$ in $\RR^2$.
\end{abstract}

%%%%%%%%%%%%%%%%%%%%%%%%%%%%%%%%%%%%%%%%%%%%%%%%%%%%%%%%%%
%%
%% MARK: Introduction
%%
%%%%%%%%%%%%%%%%%%%%%%%%%%%%%%%%%%%%%%%%%%%%%%%%%%%%%%%%%%

This question was raised by Professor Jedrzej Sniatycki:
describe the differential forms on the space $X$,
the union of the two axes in $\RR^2$
$$
  X = \{ (x,0) \mid x \in \RR\} \cup \{ (0,y) \mid y \in \RR\},
  $$
equipped with the subset diffeology.

\begin{Article}{The forms on the cross.}{}
  Every differential $1$-form $\alpha$ on $X$
  is the restriction of a $1$-form on $\RR^2$ of the type $A = a(x) dx + b(y) dy$.
  In other words,
  for all plots $P \colon r \mapsto (x_r,y_r)$ in $X$,
  $$
    \alpha(P)_r(\delta r) = a(x_r){\partial x_r \over \partial r}(\delta r) +  b(y_r){\partial y_r \over \partial r}(\delta r).
  $$
  The $1$-forms $a$ and $b$ are the restrictions of $\alpha$ to the axes:
  $$
    a(x) dx = \alpha \restriction ox \quad \text{and} \quad b(y) dy = \alpha \restriction oy.
  $$
  \Note According to the definition \cite[\textsection\,6.40]{PIZ13},
  the value of $\alpha$ at $(0,0)$ is zero if and only if $a(0)=b(0)=0$.
  Indeed, evaluated on the plots $t \mapsto (t,0)$
  and $t \mapsto (0,t)$ at $t=0$,
  $\alpha$ is equal to $a(0)dx$ or $b(0)dy$,
  which are not zero
  unless both are $0$.
\end{Article}

\begin{proof}
  First of all, notice that the four semi-axes:
  $ox^-$,
  $ox^+$,
  $oy^-$ and $oy^+$,
  are D-open in $X$.
  That is,
  open for the D-topology \cite[\textsection\,2.8]{PIZ13}.
  Indeed,
  for example,
  for any plot $P$ in $X$,
  $P^{-1}(ox^+) = P^{-1}\{(x,y) \mid x>0\}$,
  which is open because it is the pullback of an open subset by a continuous map.
  And \textit{mutatis mutandis} for $ox^-$,
  $oy^-$ and $oy^+$.
  
  Let $P$ be defined on $\U \subset \RR^n$,
  and let
  $$
    \cO = P^{-1}(ox^- \cup ox^+ \cup oy^- \cup oy^+).
  $$
  The subset $\cO$ is open.
  Let us prove first that
  $$
    \alpha - A \restriction \cO = 0.
  $$
  That is,
  for all $r \in \cO$ and all $\delta r \in \RR^n$
  $$
    \alpha(P)_r(\delta r) - a(x_r) \delta x_r - b(y_r)\delta y_r = 0.
  $$
  We denote
  $$
    \delta x_r = {\partial x_r \over \partial r}(\delta r)
    \quad \text{and} \quad
    \delta y_r = {\partial y_r \over \partial r}(\delta r).
  $$
  Let $r \in \cO$ and suppose $P(r) \in ox^+$,
  for example.
  Then,
  there exists an open neighbourhood $V$ of $r$ such that $P \restriction V$ takes its values in $ox^+$.
  Thus,
  on $V$,
  by definition $\alpha(P)_r(\delta r) = a(x_r)\delta x_r$.
  And since $y_r = 0$ on $V$,
  $\alpha(P)_r(\delta r) = a(x_r)\delta x_r + b(y_r)\delta y_r$.
  \textit{Mutatis mutandis} for all other values of $P(r)$ when $r \in \cO$.
  We eventually obtain that on $\cO$,
  $(\alpha - A)(P) \restriction \cO = 0$.
  
  Now, for all $\delta r \in \RR^n$,
  $(\alpha - A)(P)_r(\delta r)$ is a smooth function in $r$.
  This function vanishes on $\cO$;
  by continuity it vanishes on the closure $\overline\cO$.
  Let $\U' = \U \setminus \overline\cO$,
  which is an open subset of $\RR^n$.
  Then $P \restriction \U'$ is a plot in $X$ taking only the value $(0,0)$.
  Hence,
  for all $r \in U'$,
  $(\alpha - A)(P)_r(\delta r) = 0$,
  since a differential form evaluated on a constant plot is zero (as it factors through $\RR^0$).
  Thus,
  for all $r \in U$,
  $(\alpha - A)(P)_r(\delta r) = 0$.
  Hence $\alpha(P)_r(\delta r) = a(x_r) \delta x_r + b(y_r) \delta y_r$,
  and $\alpha = A \restriction X$.
\end{proof}

\vfill
\includegraphics[width=\textwidth]{Figures/Thecross.pdf}
